\PassOptionsToPackage{table}{xcolor}
\documentclass[10pt,aspectratio=169]{beamer}
\newcommand{\LectureNumber}{19}
\input{course-theme.tex}
\title{Recursion and backtracking}
\subtitle{CSC103 Programming Fundamentals / Lecture 19}
\author{Dr. Muhammad Farhan}
\institute{Associate Professor of Computer Science\\Department of Computer Science\\COMSATS University Islamabad\\Sahiwal Campus}
\date{Fall 2026}
\begin{document}
\begin{frame}\titlepage\end{frame}

\begin{frame}[fragile,t]{The problem for this lecture}
\begin{columns}[T,onlytextwidth]\begin{column}{0.60\textwidth}
Trace factorial(4) with base case factorial(0)=1. Then write the recursive method for inputs 0..12 using int.
\end{column}\begin{column}{0.35\textwidth}\small
Reject n outside the supported range before computing.
\end{column}\end{columns}
\note{Introduce the target problem before explaining the tools. Revisit it in the guided solution later.\par Syllabus reading: Deitel Ch 18. CLO-2.}
\end{frame}

\begin{frame}[fragile,t]{Learning objectives}
\begin{columns}[T,onlytextwidth]\begin{column}{0.60\textwidth}
\begin{itemize}
\item Identify base and recursive cases
\item Trace recursive stack frames
\item Explain recursive search and backtracking
\end{itemize}
\end{column}\begin{column}{0.35\textwidth}\small
CLO-2

Base case and progress\par Recursive call stack\par Recursion and iteration\par Backtracking search
\end{column}\end{columns}
\note{Explain how each objective supports the opening problem. The final questions check these objectives.\par Syllabus reading: Deitel Ch 18. CLO-2.}
\end{frame}

\begin{frame}[fragile,t]{Base case and progress}
\begin{itemize}
\item A recursive method calls itself directly or indirectly.
\item A base case returns without another recursive call.
\item Every permitted input must progress toward a base case.
\end{itemize}
\vspace{1mm}\begin{block}{Example}
\begin{lstlisting}
sumTo(0) = 0
sumTo(n) = n + sumTo(n-1), n > 0
\end{lstlisting}
\end{block}
\note{Use nonnegative n. Without a domain restriction, decrementing a negative argument moves away from zero.\par Syllabus reading: Deitel Ch 18. CLO-2.}
\end{frame}

\begin{frame}[fragile,t]{Recursive call stack}
\begin{itemize}
\item Each recursive call keeps its own parameters and pending work.
\item Returns unwind in reverse order.
\item The deepest base case determines the first returned result.
\end{itemize}
\vspace{1mm}\begin{block}{Example}
\begin{lstlisting}
sumTo(3): waits for 3 + sumTo(2)
sumTo(2): waits for 2 + sumTo(1)
sumTo(1): waits for 1 + sumTo(0)
sumTo(0): returns 0
\end{lstlisting}
\end{block}
\note{Then return 1,3,6. Ask students to identify which frame owns each n rather than imagining a single shared n.\par Syllabus reading: Deitel Ch 18. CLO-2.}
\end{frame}

\begin{frame}[fragile,t]{Recursion and iteration}
\begin{itemize}
\item Both can express repetition.
\item A loop can often use constant auxiliary space for a simple sum.
\item Recursive depth consumes stack space and can cause StackOverflowError.
\end{itemize}
\vspace{1mm}\begin{block}{Example}
\begin{lstlisting}
Iterative sum:
int total = 0;
for (int i=1; i<=n; i++) total += i;
\end{lstlisting}
\end{block}
\note{For this direct recursive sum, depth grows with n. Java does not guarantee tail-call optimisation. Use small demonstration inputs.\par Syllabus reading: Deitel Ch 18. CLO-2.}
\end{frame}

\begin{frame}[fragile,t]{Backtracking search}
\begin{itemize}
\item Backtracking explores a choice, then explores alternatives when needed.
\item A base case recognises a complete candidate.
\item Immutable partial results naturally preserve the earlier choice state.
\end{itemize}
\vspace{1mm}\begin{block}{Example}
\begin{lstlisting}
binary("", 2)
  binary("0", 1)
    outputs 00, 01
  binary("1", 1)
    outputs 10, 11
\end{lstlisting}
\end{block}
\note{Algorithm: if remaining==0 print prefix; otherwise call binary(prefix+"0",remaining-1), then prefix+"1". Mutable state would need explicit undo. This is exhaustive branching search.\par Syllabus reading: Deitel Ch 18. CLO-2.}
\end{frame}

\begin{frame}[fragile,t]{A recursive definition needs a domain}
\begin{itemize}
\item factorial(0)=1 supplies the base case.
\item For positive n, factorial(n)=n*factorial(n-1).
\item For this int implementation, restrict n to 0..12 so that results fit the chosen type.
\end{itemize}
\vspace{1mm}\begin{block}{Example}
\begin{lstlisting}
factorial(4) = 4 * 3 * 2 * 1 = 24
13! exceeds the positive int range.
\end{lstlisting}
\end{block}
\note{Explain the mechanism using the displayed example. factorial(0)=1 supplies the base case. For positive n, factorial(n)=n*factorial(n-1). For this int implementation, restrict n to 0..12 so that results fit the chosen type.\par Syllabus reading: Deitel Ch 18. CLO-2.}
\end{frame}

\begin{frame}[fragile,t]{Returning is different from making a call}
\begin{itemize}
\item A factorial frame waits for the smaller factorial before multiplying.
\item The base case provides the first completed answer.
\item Unwinding performs the pending multiplications from the deepest frame upward.
\end{itemize}
\vspace{1mm}\begin{block}{Example}
\begin{lstlisting}
fact(0) returns 1
fact(1) returns 1 * 1
fact(2) returns 2 * 1
fact(3) returns 3 * 2
\end{lstlisting}
\end{block}
\note{Explain the mechanism using the displayed example. A factorial frame waits for the smaller factorial before multiplying. The base case provides the first completed answer. Unwinding performs the pending multiplications from the deepest frame upward.\par Syllabus reading: Deitel Ch 18. CLO-2.}
\end{frame}


\begin{frame}[fragile,t]{Recursive calls unwind in reverse order}
\begin{center}\begin{tikzpicture}
\node[box,text width=60mm] (n0) at (0,0.0) {fact(3)};
\node[box,text width=60mm] (n1) at (0,-1.05) {fact(2)};
\node[box,text width=60mm] (n2) at (0,-2.1) {fact(1)};
\node[alt,text width=60mm] (n3) at (0,-3.1500000000000004) {fact(0) = 1};
\draw[arr] (n0.east)--++(1,0)|-(n1.east);
\draw[arr] (n1.east)--++(1,0)|-(n2.east);
\draw[arr] (n2.east)--++(1,0)|-(n3.east);
\draw[arr] (n3.west)--++(-1,0)|-node[lab,pos=.5]{returns unwind}(n0.west);
\end{tikzpicture}\end{center}
\begin{block}{What to notice}Calls move toward the base case; returned values then flow back to waiting callers.\end{block}
\note{Trace each labelled connection. Calls move toward the base case; returned values then flow back to waiting callers.}
\end{frame}
\begin{frame}[fragile,t]{Key distinctions}
\small\renewcommand{\arraystretch}{1.5}
\rowcolors{2}{pale}{white}
\begin{tabularx}{\textwidth}{>{\raggedright\arraybackslash}X>{\raggedright\arraybackslash}X>{\raggedright\arraybackslash}X}
\rowcolor{navy}
\color{white}\bfseries Technique & \color{white}\bfseries Stored work & \color{white}\bfseries Good fit\\
Iteration & Loop state and accumulator & Simple repeated accumulation\\
Direct recursion & Pending work in frames & A problem with a smaller same-form case\\
Backtracking & Choices and search state & Exploring alternative candidates\\
\end{tabularx}
\vspace{3mm}\footnotesize These distinctions determine which operation or control structure fits the problem.
\note{\par Syllabus reading: Deitel Ch 18. CLO-2.}
\end{frame}

\begin{frame}[fragile,t]{First example: methods}
\begin{lstlisting}
static int sumTo(int n) {
  if (n < 0) throw new IllegalArgumentException("negative n");
  if (n == 0) return 0;
  return n + sumTo(n - 1);
}
static void binary(String prefix, int remaining) {
  if (remaining == 0) { System.out.println(prefix); return; }
  binary(prefix + "0", remaining - 1);
  binary(prefix + "1", remaining - 1);
}
\end{lstlisting}
\note{Method definitions belong inside the class and outside main. Explain the parameters and return values.\par Syllabus reading: Deitel Ch 18. CLO-2.}
\end{frame}

\begin{frame}[fragile,t]{First example: execution}
\begin{lstlisting}
System.out.println(sumTo(4));
binary("", 2);
\end{lstlisting}
\note{This is the main body from Java\_Examples/Lecture19.java. The source file includes the complete class. Predict the result before the next slide.\par Syllabus reading: Deitel Ch 18. CLO-2.}
\end{frame}

\begin{frame}[fragile,t]{First example: result}
\begin{columns}[T,onlytextwidth]\begin{column}{0.60\textwidth}
10\par 00\par 01\par 10\par 11\par Each binary choice reduces the remaining depth by one.
\end{column}\begin{column}{0.35\textwidth}\small
Recursion and iteration

Backtracking search
\end{column}\end{columns}
\note{Expected output and interpretation: 10
00
01
10
11
Each binary choice reduces the remaining depth by one.\par Syllabus reading: Deitel Ch 18. CLO-2.}
\end{frame}

\begin{frame}[fragile,t]{Guided practice}
\begin{columns}[T,onlytextwidth]\begin{column}{0.60\textwidth}
Trace factorial(4) with base case factorial(0)=1. Then write the recursive method for inputs 0..12 using int.
\end{column}\begin{column}{0.35\textwidth}\small
Attempt the solution before continuing.

Use the definitions and examples from this lecture.
\end{column}\end{columns}
\note{Give students 8 minutes to attempt the problem. The following slides provide a complete solution. Original solution guidance: Pending multiplications: 4*fact(3), 3*fact(2), 2*fact(1), 1*fact(0). Returns: 1,1,2,6,24. Larger values can overflow int.\par Syllabus reading: Deitel Ch 18. CLO-2.}
\end{frame}

\begin{frame}[fragile,t]{Solution strategy}
\begin{columns}[T,onlytextwidth]\begin{column}{0.60\textwidth}
\begin{itemize}
\item Reject n outside the supported range before computing.
\item Return 1 at n==0.
\item For a positive n, return n multiplied by factorial(n-1).
\end{itemize}
\end{column}\begin{column}{0.35\textwidth}\small
The next program uses fixed inputs so its result can be reproduced.
\end{column}\end{columns}
\note{Discuss why each step is needed. State the input assumptions before executing the program.\par Syllabus reading: Deitel Ch 18. CLO-2.}
\end{frame}

\begin{frame}[fragile,t]{Complete solution}
\begin{lstlisting}
public class Solution19 {
  public static void main(String[] args) throws Exception {
    System.out.println(factorial(4));
    System.out.println(factorial(0));
  }
  static int factorial(int n) {
    if (n < 0 || n > 12) {
      throw new IllegalArgumentException("n must be in 0..12");
    }
    if (n == 0) return 1;
    return n * factorial(n - 1);
  }
}
\end{lstlisting}
\note{Complete runnable file: Java\_Examples/Solution19.java. Inputs are fixed in the program.  \par Syllabus reading: Deitel Ch 18. CLO-2.}
\end{frame}

\begin{frame}[fragile,t]{Execution trace}
\small\renewcommand{\arraystretch}{1.5}
\rowcolors{2}{pale}{white}
\begin{tabularx}{\textwidth}{>{\raggedright\arraybackslash}X>{\raggedright\arraybackslash}X>{\raggedright\arraybackslash}X}
\rowcolor{navy}
\color{white}\bfseries Returning frame & \color{white}\bfseries Calculation & \color{white}\bfseries Returned value\\
factorial(0) & Base case & 1\\
factorial(1) & 1 * 1 & 1\\
factorial(2) & 2 * 1 & 2\\
factorial(3) & 3 * 2 & 6\\
factorial(4) & 4 * 6 & 24\\
\end{tabularx}
\vspace{3mm}\footnotesize Follow the changing state in execution order.
\note{Manually derived trace for the complete solution. Reject n outside the supported range before computing. Return 1 at n==0. For a positive n, return n multiplied by factorial(n-1).\par Syllabus reading: Deitel Ch 18. CLO-2.}
\end{frame}

\begin{frame}[fragile,t]{Solution result}
\begin{columns}[T,onlytextwidth]\begin{column}{0.60\textwidth}
24\par 1
\end{column}\begin{column}{0.35\textwidth}\small
Why this result follows

For a positive n, return n multiplied by factorial(n-1).
\end{column}\end{columns}
\note{Expected result: 24
1\par Syllabus reading: Deitel Ch 18. CLO-2.}
\end{frame}

\begin{frame}[fragile,t]{A common incorrect approach}
static int f(int n) \{ return n + f(n-1); \}
\vspace{1mm}\begin{block}{Example}
\begin{lstlisting}
There is no base case. Add an appropriate base case and enforce a domain that progresses toward it.
\end{lstlisting}
\end{block}
\note{Ask students to predict the failure before discussing the explanation. The right side gives the correction.\par Syllabus reading: Deitel Ch 18. CLO-2.}
\end{frame}

\begin{frame}[fragile,t]{Boundary and failure checks}
\begin{columns}[T,onlytextwidth]\begin{column}{0.60\textwidth}
Check n=0, n=1 and n=4. Then verify rejection for n=-1 and n=13.
\end{column}\begin{column}{0.35\textwidth}\small
Pending multiplications: 4*fact(3), 3*fact(2), 2*fact(1), 1*fact(0). Returns: 1,1,2,6,24. Larger values can overflow int.
\end{column}\end{columns}
\note{Use the specification to derive each expected result. Exercise solution: Pending multiplications: 4*fact(3), 3*fact(2), 2*fact(1), 1*fact(0). Returns: 1,1,2,6,24. Larger values can overflow int.\par Syllabus reading: Deitel Ch 18. CLO-2.}
\end{frame}

\begin{frame}[fragile,t]{Check your understanding}
\begin{columns}[T,onlytextwidth]\begin{column}{0.60\textwidth}
Why does backtracking need to restore mutable state? What causes the direct recursive sum to terminate?
\end{column}\begin{column}{0.35\textwidth}\small
Write a prediction and explain the rule that supports it.
\end{column}\end{columns}
\note{Answer: A later branch must start from the state before the earlier choice. sumTo reduces nonnegative n until n==0.\par Syllabus reading: Deitel Ch 18. CLO-2.}
\end{frame}

\begin{frame}[fragile,t]{Answer and explanation}
\begin{columns}[T,onlytextwidth]\begin{column}{0.60\textwidth}
A later branch must start from the state before the earlier choice. sumTo reduces nonnegative n until n==0.
\end{column}\begin{column}{0.35\textwidth}\small
There is no base case. Add an appropriate base case and enforce a domain that progresses toward it.
\end{column}\end{columns}
\note{Reconnect the answer to the relevant definition rather than treating it as a fact to memorise.\par Syllabus reading: Deitel Ch 18. CLO-2.}
\end{frame}


\begin{frame}[fragile,t]{Compare cases and predict outcomes}
\small\renewcommand{\arraystretch}{1.5}\rowcolors{2}{pale}{white}
\begin{tabularx}{\textwidth}{>{\raggedright\arraybackslash}X>{\raggedright\arraybackslash}X>{\raggedright\arraybackslash}X}\rowcolor{navy}
\color{white}\bfseries Returning from & \color{white}\bfseries Calculation & \color{white}\bfseries Returned value\\
fact(1) & 1 * 1 & 1\\
fact(2) & 2 * 1 & 2\\
fact(3) & 3 * 2 & 6\\
\end{tabularx}\par\vspace{3mm}\begin{exampleblock}{Discuss}Explain each expected result before running the example. Which case would reveal a wrong assumption?\end{exampleblock}
\note{Read the first two columns and invite predictions before discussing the outcome.}
\end{frame}

\begin{frame}[fragile,t]{Apply the idea}
\begin{alertblock}{Independent practice}Write recursive sumTo(n) for 0 through 100. Trace sumTo(3), including the base-case return.\end{alertblock}\vspace{3mm}\begin{enumerate}\item Work through the input by hand.\item Write or correct the program.\item Justify the expected result.\end{enumerate}\note{Allow 3–5 minutes, then compare answers in pairs. Write recursive sumTo(n) for 0 through 100. Trace sumTo(3), including the base-case return.}
\end{frame}

\begin{frame}[fragile,t]{Solution and reasoning}
\begin{lstlisting}
static int sumTo(int n) {
  if (n < 0 || n > 100)
    throw new IllegalArgumentException("n outside 0..100");
  if (n == 0) return 0;
  return n + sumTo(n - 1);
}
\end{lstlisting}\vspace{1mm}\small\begin{itemize}
\item The calls are sumTo(3), sumTo(2), sumTo(1), sumTo(0).
\item The returned values are 0, 1, 3 and 6.
\item The stated bound limits recursion depth and keeps these sums within int.
\end{itemize}\note{Answer: The calls are sumTo(3), sumTo(2), sumTo(1), sumTo(0). The returned values are 0, 1, 3 and 6. The stated bound limits recursion depth and keeps these sums within int.}
\end{frame}

\begin{frame}[fragile,t]{Fibonacci recursion and repeated work}
\begin{itemize}
\item Two base cases define fib(0)=0 and fib(1)=1.
\item For n\textgreater{}=2, two smaller calls contribute to the result.
\item Naive recursion repeats subproblems; use only small n for the demonstration.
\end{itemize}
\vspace{3mm}\footnotesize\textcolor{accent}{Source: supplied ISB campus slides. See speaker notes and ISB\_Source\_Map.md.}
\note{Adapted from the supplied ISB campus folder: Recursion(New).pptx, slides 33, 34, 35; Java Recursion.pptx, slides 6, 7, 12, 13. Added a small demonstration bound. Memoization and tree traversal in the source are optional enrichment, not new assessed prerequisites.}
\end{frame}

\begin{frame}[fragile,t]{Worked problem: Fibonacci recursion and repeated work}
\begin{block}{Task}Calculate fib(6), and trace fib(4) to see fib(2) computed more than once.\end{block}
\small Input: Values are fixed in the program for a reproducible demonstration.\par\vspace{3mm}\begin{exampleblock}{Before running}Predict the output and identify the variables that change.\end{exampleblock}
\note{Adapted from the supplied ISB campus folder: Recursion(New).pptx, slides 33, 34, 35; Java Recursion.pptx, slides 6, 7, 12, 13. Added a small demonstration bound. Memoization and tree traversal in the source are optional enrichment, not new assessed prerequisites.}
\end{frame}

\begin{frame}[fragile,t]{Complete ISB example (1/1)}
\begin{lstlisting}
public class IsbLecture19 {
  public static void main(String[] args) throws Exception {
    System.out.println(fib(6));
  }
  static int fib(int n) {
    if (n < 0 || n > 30) throw new IllegalArgumentException("use 0..30");
    if (n < 2) return n;
    return fib(n - 1) + fib(n - 2);
  }
}
\end{lstlisting}
\note{Adapted from the supplied ISB campus folder: Recursion(New).pptx, slides 33, 34, 35; Java Recursion.pptx, slides 6, 7, 12, 13. Added a small demonstration bound. Memoization and tree traversal in the source are optional enrichment, not new assessed prerequisites. Complete file: ISB\_Java\_Examples/IsbLecture19.java.}
\end{frame}

\begin{frame}[fragile,t]{Worked trace and expected result}
\small\renewcommand{\arraystretch}{1.4}\rowcolors{2}{pale}{white}
\begin{tabularx}{\textwidth}{>{\raggedright\arraybackslash}X>{\raggedright\arraybackslash}X>{\raggedright\arraybackslash}X}\rowcolor{navy}
\color{white}\bfseries Index & \color{white}\bfseries Recurrence / base & \color{white}\bfseries Value\\
0 / 1 & Base cases & 0 / 1\\
2 / 3 & 1+0 / 1+1 & 1 / 2\\
4 / 5 / 6 & 2+1 / 3+2 / 5+3 & 3 / 5 / 8\\
\end{tabularx}\par\vspace{2mm}
\begin{block}{Expected console output}\ttfamily\footnotesize 8\end{block}
\note{Adapted from the supplied ISB campus folder: Recursion(New).pptx, slides 33, 34, 35; Java Recursion.pptx, slides 6, 7, 12, 13. Added a small demonstration bound. Memoization and tree traversal in the source are optional enrichment, not new assessed prerequisites. Trace and expected values independently worked for this edition.}
\end{frame}

\begin{frame}[fragile,t]{Extend the ISB example}
\begin{alertblock}{Practice}Why does changing the method to tail{-}recursive form not guarantee constant stack space in Java?\end{alertblock}\vspace{3mm}\begin{exampleblock}{Explain your reasoning}Java does not guarantee tail{-}call elimination. Use an explicit loop when bounded stack usage is required.\end{exampleblock}
\note{Adapted from the supplied ISB campus folder: Recursion(New).pptx, slides 33, 34, 35; Java Recursion.pptx, slides 6, 7, 12, 13. Added a small demonstration bound. Memoization and tree traversal in the source are optional enrichment, not new assessed prerequisites. Answer: Java does not guarantee tail{-}call elimination. Use an explicit loop when bounded stack usage is required.}
\end{frame}
\begin{frame}[fragile,t]{Lecture summary}
\begin{columns}[T,onlytextwidth]\begin{column}{0.60\textwidth}
\begin{itemize}
\item base and recursive cases
\item recursive stack frames
\item recursive search and backtracking
\end{itemize}
\end{column}\begin{column}{0.35\textwidth}\small
\begin{itemize}
\item Reject n outside the supported range before computing.
\item Return 1 at n==0.
\item For a positive n, return n multiplied by factorial(n-1).
\end{itemize}
\end{column}\end{columns}
\note{Ask students to give one concrete example for the concepts listed. Use the worked solution to connect the topics.\par Syllabus reading: Deitel Ch 18. CLO-2.}
\end{frame}

\begin{frame}[fragile,t]{After-class practice}
\begin{columns}[T,onlytextwidth]\begin{column}{0.60\textwidth}
Generate all three-character binary strings recursively. Explain why there are eight results.
\end{column}\begin{column}{0.35\textwidth}\small
Syllabus reading\par Deitel Ch 18

Next lecture: One-dimensional arrays
\end{column}\end{columns}
\note{Reading labels follow the supplied outline. Check topic headings against the available textbook edition. Require a program and expected results for the practice task.\par Syllabus reading: Deitel Ch 18. CLO-2.}
\end{frame}
\end{document}
